% Preliminaries: the structural lemmas for a(x-y) = bz.
% Written by the orchestrator, 2026-08-12. Every statement here is proved;
% the computational checks recorded alongside each are consistency tests,
% not the justification.
%
% NOVELTY STATUS, stated plainly so it cannot drift in editing:
%   Lemma  (solution form)   -- elementary, certainly folklore.
%   Cor    (a | z)           -- immediate from the above.
%   Lemma  (self-similarity) -- elementary.
%   Prop   (lifting)         -- a proof technique, not a new bound.
%   Cor    (R_k >= a^k)      -- THIS IS CHANG-DE LOERA-WESLEY LEMMA 4.1.
%                               Reproved here, NOT new. Do not present it
%                               as a contribution of this work.

\section{Preliminaries}

Fix integers $a \ge 2$ and $b \ge 1$ with $\gcd(a,b) = 1$, and consider the
equation
\[
  \mathcal{E}(a,b): \qquad a(x - y) = b z .
\]
A \emph{solution in $[n]$} is a triple $(x,y,z) \in [n]^3$ satisfying
$\mathcal{E}(a,b)$; the entries need not be distinct, and in particular $z$
may coincide with $x$ or $y$. Given a colouring $\chi : [n] \to [k]$, a
solution is \emph{monochromatic} if $\chi(x) = \chi(y) = \chi(z)$, and $\chi$
is \emph{solution-free} if it admits none. The \emph{$k$-colour Rado number}
$R_k(\mathcal{E})$ is the least $n$ such that every $\chi : [n] \to [k]$
admits a monochromatic solution; exhibiting a solution-free $k$-colouring of
$[N]$ therefore certifies $R_k(\mathcal{E}) > N$.

Write $v(j)$ for the $a$-adic valuation of $j$, that is, the largest $e$ with
$a^e \mid j$, and call $j$ a \emph{unit} when $v(j) = 0$.

\subsection{The shape of a solution}

\begin{lemma}[Solution form]\label{lem:solform}
Let $\gcd(a,b) = 1$. A triple $(x,y,z)$ of positive integers satisfies
$a(x-y) = bz$ if and only if there is an integer $t \ge 1$ with
\[
  x - y = b t, \qquad z = a t .
\]
\end{lemma}

\begin{proof}
Suppose $a(x-y) = bz$. Then $a \mid bz$, and since $\gcd(a,b) = 1$ we get
$a \mid z$; write $z = at$. Substituting gives $a(x-y) = bat$, and cancelling
$a \neq 0$ leaves $x - y = bt$. As $z \ge 1$ and $a \ge 2$ we have $t \ge 1$.
The converse is a direct substitution:
$a(bt) = b(at)$.
\end{proof}

Lemma~\ref{lem:solform} has an immediate consequence which, despite its
triviality, carries most of the weight in what follows.

\begin{corollary}\label{cor:adivz}
Every solution of $\mathcal{E}(a,b)$ in the positive integers has $a \mid z$.
Equivalently, \emph{a unit is never the $z$-component of a solution.}
\end{corollary}

Consequently a colour class consisting entirely of units is automatically
free of monochromatic solutions: a monochromatic solution in that class would
supply a unit $z$, contradicting Corollary~\ref{cor:adivz}. This observation
is what makes the valuation-based constructions below work, and it is the
reason a case analysis over colour classes need only treat classes meeting
$a\mathbb{Z}$.

\begin{lemma}[Self-similarity]\label{lem:selfsim}
Let $(x,y,z)$ be a solution of $\mathcal{E}(a,b)$ with $a \mid x$, $a \mid y$
and $a \mid z$. Then $(x/a,\, y/a,\, z/a)$ is again a solution of
$\mathcal{E}(a,b)$.
\end{lemma}

\begin{proof}
By Lemma~\ref{lem:solform} write $x - y = bt$ and $z = at$. Since $a \mid x$
and $a \mid y$ we have $a \mid bt$, and $\gcd(a,b) = 1$ forces $a \mid t$;
write $t = a t'$ with $t' \ge 1$. Then
$x/a - y/a = (x-y)/a = bt' $ and $z/a = t = a t'$, which is precisely the
form of Lemma~\ref{lem:solform} with parameter $t'$.
\end{proof}

Thus the multiples of $a$ inside $[n]$ carry a scaled copy of the same
problem on $[\lfloor n/a \rfloor]$. This self-similarity is the engine of the
following lifting principle.

\subsection{A lifting principle}

\begin{proposition}[Lifting]\label{prop:lift}
Suppose $\chi'$ is a solution-free $k$-colouring of $[M]$. Define a
$(k+1)$-colouring $\chi$ of $[\,aM + a - 1\,]$ by
\[
  \chi(j) =
  \begin{cases}
    \star & \text{if } a \nmid j, \\[2pt]
    \chi'(j/a) & \text{if } a \mid j,
  \end{cases}
\]
where $\star$ is a colour not used by $\chi'$. Then $\chi$ is solution-free.
\end{proposition}

\begin{proof}
Suppose $(x,y,z)$ were a monochromatic solution in $[\,aM+a-1\,]$, of common
colour $c$.

If $c = \star$ then $a \nmid z$, contradicting Corollary~\ref{cor:adivz}.

Otherwise $c \neq \star$, so none of $x,y,z$ receives $\star$ and hence
$a$ divides each of them. By Lemma~\ref{lem:selfsim} the triple
$(x/a, y/a, z/a)$ is a solution. Each component is a multiple of $a$ bounded
by $aM + a - 1$, hence bounded by $aM$, so each of $x/a, y/a, z/a$ lies in
$[M]$. Finally
$\chi'(x/a) = \chi(x) = c$ and likewise for $y$ and $z$, so $\chi'$ admits a
monochromatic solution in $[M]$ --- a contradiction.
\end{proof}

\begin{lemma}[Base case]\label{lem:base}
A single colour on $[N]$ is solution-free if and only if
$N \le \max(a-1,\, b)$.
\end{lemma}

\begin{proof}
By Lemma~\ref{lem:solform} the solution with the smallest components is the
one with $t = 1$, namely $x = y + b$, $z = a$ with $y \ge 1$. It lies in
$[N]$ exactly when $N \ge a$ and $N \ge b+1$. Hence a solution exists in
$[N]$ if and only if $N \ge \max(a, b+1)$, and $[N]$ is solution-free if and
only if $N \le \max(a, b+1) - 1 = \max(a-1, b)$.
\end{proof}

\begin{corollary}\label{cor:ak}
Let $2 \le a$ and $1 \le b \le a-1$ with $\gcd(a,b) = 1$. Then for every
$k \ge 1$,
\[
  R_k\bigl(a(x-y) = bz\bigr) \;\ge\; a^k .
\]
\end{corollary}

\begin{proof}
Let $f(k)$ denote the bound obtained by iterating
Proposition~\ref{prop:lift} from Lemma~\ref{lem:base}. Since $b \le a-1$ we
have $\max(a-1,b) = a-1$, so $f(1) = a - 1$, and
Proposition~\ref{prop:lift} gives $f(k+1) = a f(k) + a - 1$. By induction
$f(k) = a^k - 1$: it holds at $k=1$, and
$a(a^k - 1) + a - 1 = a^{k+1} - 1$. Unwinding the lifting, the resulting
colouring is $\chi(j) = \min(v(j),\, k-1)$. It is a solution-free
$k$-colouring of $[a^k - 1]$, whence $R_k > a^k - 1$.
\end{proof}

\begin{remark}\label{rem:notnew}
Corollary~\ref{cor:ak} is \emph{not} new: it is Lemma~4.1 of
Chang, De~Loera and Wesley~\cite{CDW2022}. We have reproved it here because
the inductive argument, and especially Corollary~\ref{cor:adivz}, are the
tools the construction of Section~\ref{sec:construction} builds on. We note
also that the bound is exactly tight for this particular colouring: a direct
check confirms that $\chi(j) = \min(v(j), k-1)$ is solution-free on
$[a^k-1]$ and \emph{fails} on $[a^k]$, for every $a \le 6$, $b < a$ with
$\gcd(a,b)=1$, and $k \le 5$. Any improvement must therefore come from a
different colouring, not a sharper analysis of this one.
\end{remark}
